% !PRINTED none
\weberpage{none}
% [source note: this Abschnitt opener has no printed folio; sequence position corresponds to inferred page 705.]

\begin{center}
\gesperrt{Twenty-fourth Chapter.}

\gesperrt{Class number of the cyclotomic fields.}

\rule{4em}{0.4pt}
\end{center}

\begin{center}
\S.\,193.

\gesperrt{Class number representation in the cyclotomic field $\Omega_m$.}
\end{center}

The general results concerning the class number $h$ are now to
be applied to the full cyclotomic field $\Omega_m$, under the
assumption made hitherto that $m$ is a power of a prime
$q$, and that accordingly
\weq{1}{m = q^{\varkappa}, \quad \varphi(m) = q^{\varkappa-1}(q-1) = \mu}
has been set.

In the twenty-first chapter (\S.\,169, 171) we have seen
that only a single prime ideal of degree $1$ divides $q$, and
that a prime $p$ \gesperrt{belonging} to the exponent $f$, and
different from $q$, splits into $e$ distinct prime ideals of degree $f$.
Here $f$ denotes the least positive exponent for which
\weq{2}{p^f \equiv 1\ (\mathrm{mod}\ m)}
holds, and the equation
\weq{3}{\mu = ef}
determines $e$.

The function $\varPhi(s)$ determined at the end of the preceding
section by formula (6) therefore takes here the expression:
\weq{4}{\varPhi(s) = \frac{1}{1-q^{-s}}\prod\frac{1}{(1-p^{-sf})^e},}
and in it the product sign $\varPi$ extends over all primes
 $p$ that are different from $q$. Here $s$ is a variable
that is always greater than 1, and that is finally to approach
the limit 1.
